\documentclass[11pt,a4paper]{article}
\usepackage[margin=25mm]{geometry}
\usepackage[T1]{fontenc}
\usepackage{lmodern}
\usepackage{amsmath,amssymb}
\usepackage{hyperref}
\hypersetup{hidelinks,
  pdftitle={Machine Learning for Kira IBP Reduction},
  pdfsubject={Feynman integrals, IBP reduction, and proposed machine learning roles}}
\setlength{\parindent}{0pt}
\setlength{\parskip}{6pt}
\setlength{\emergencystretch}{2em}
\title{\textbf{Machine Learning for Kira IBP Reduction}\\
  \large Seed selection and elimination ordering}
\author{}
\date{October 1, 2026}

\begin{document}
\maketitle
\vspace{-1em}

\begin{abstract}
Feynman integral calculations are an important component of multiloop amplitude calculations in quantum field theory. In Quantum Chromodynamics (QCD), multiloop amplitudes and their associated Feynman integrals arise in higher-order perturbative calculations and can become increasingly complex at higher loop orders. Recent progress has automated these calculations to some extent using the framework Kira, which organizes target integrals into integral families and sectors, generate algebraic relations among the integrals by using integration by parts (IBP) identities, and reduce the resulting system to expressions in terms of a finite set of master integrals. We aim to explore the application of System 1 machine learning models to accelerate the IBP reduction process by investigating if lightweight learned models can improve seed selection. These models can predict the most useful integrals from which to generate IBP relations while minimizing redundant calculations. We will also explore System 1 models for learning efficient reduction and elimination sequences to minimize the number of algebraic operations required to express target integrals in terms of master integrals. Because the machine learning models will serve as decision-makers while the underlying IBP reduction and verification remain exact, these models can improve computational efficiency without compromising the correctness of the resulting reductions.
\end{abstract}

\clearpage
\section{Introduction}
\label{sec:introduction}
Feynman integrals arise in higher-order perturbative corrections to
scattering amplitudes and particle decay amplitudes. In QCD, these
corrections involve quark and gluon loops and contribute to predictions
for cross sections and decay rates. Calculating a loop contribution
requires integration over the independent internal momenta that remain
after imposing momentum conservation~\cite{Weinzierl2010}.

The scalar integrals that appear in an amplitude have the general form
\begin{equation}
 \mathcal J=
 \left(\prod_{k=1}^{L}\int
       \frac{\mathrm{d}^{d}\ell_k}{\mathrm{i}\pi^{d/2}}\right)
 \frac{\mathcal N(\boldsymbol\ell,\boldsymbol p)}
      {\prod_{i=1}^{n_{\mathrm{phys}}}D_i^{a_i}},
 \qquad D_i=q_i^2-m_i^2+\mathrm{i}0.
 \label{eq:general-integral}
\end{equation}
Here $L$ is the number of loops and
$\boldsymbol\ell=(\ell_1,\ldots,\ell_L)$ is the list of independent loop
momenta over which we integrate. The number of independent external
momenta is $E$, and we keep these momenta,
$\boldsymbol p=(p_1,\ldots,p_E)$, fixed during integration. The numerator
$\mathcal N$ is a polynomial obtained from the spin and tensor algebra;
we leave its mass dependence implicit in the notation.
The number of physical denominators is $n_{\mathrm{phys}}$, and each
$D_i$ corresponds to an internal particle line with momentum $q_i$ and
mass $m_i$. The positive integer $a_i$ is the power of $D_i$ in the
denominator, and $q_i$ is a fixed linear combination of loop and external
momenta, with $q_i^2=q_i\cdot q_i$ denoting the Lorentz scalar product.

The integration dimension is $d=4-2\epsilon$, where $\epsilon$ is the
regulator used to define divergent integrals, so the result depends on
$\epsilon$ as well as the external momenta and masses. The factor
$1/(\mathrm{i}\pi^{d/2})$ is our normalization convention for each
integration measure, with $\mathrm{i}^2=-1$. In each denominator,
$\mathrm{i}0$ is an infinitesimal positive imaginary term that specifies
how to integrate past a propagator pole, where $q_i^2=m_i^2$.

Higher loop orders produce collections of integrals with related numerators
and denominators and algebraic identities can be used to express them in
terms of an independent set of master integrals. Kira automates this
reduction~\cite{KiraOriginal}, after which the calculation requires the
values of the master integrals.
To explain how the reduction works and where machine learning (ML) could
speed it up, we first define the notation used in this paper, then derive
the identities and describe how Kira solves the resulting equations.

\section{The mathematical form of Feynman integrals}
\label{sec:integral-form}
To label the integrals related to Eq.~\eqref{eq:general-integral}, we
specify their momentum dependence and write the numerator and denominator
powers as a vector of integers. This vector identifies each integral as
an unknown in the reduction equations.

\subsection{Loop momenta, propagators and regularization}
After choosing how the loop and external momenta flow through the internal
lines, we write each physical denominator as
\begin{equation}
 D_i=q_i^2-m_i^2+\mathrm{i}0,
 \qquad
 q_i=\sum_{k=1}^{L}\alpha_{ik}\ell_k
       +\sum_{\beta=1}^{E}\eta_{i\beta}p_\beta .
 \label{eq:propagator}
\end{equation}
The constants $\alpha_{ik}$ and $\eta_{i\beta}$ specify which loop and
external momenta make up $q_i$, with indices $k,j$ for loop momenta and
$\beta,\gamma$ for external momenta. Squaring $q_i$ gives the scalar
products $\ell_k\cdot\ell_j$, $\ell_k\cdot p_\beta$ and
$p_\beta\cdot p_\gamma$, and those containing loop momenta vary during
integration. The vector $\boldsymbol x$ contains the independent external
scalar products and squared masses, which are the kinematic parameters
we keep fixed during integration.

Loop integrals can diverge at large momenta or when massless momenta become
small or parallel. Dimensional regularization defines these integrals at
values of $d$ for which they converge and extends the result by analytic
continuation to $d=4-2\epsilon$, using additional regulators where
necessary. In an expansion around $\epsilon=0$, the divergences appear
as negative powers of $\epsilon$, which
renormalization and infrared cancellation in the observable address later
in the calculation~\cite{Weinzierl2010}. Having specified the regulated
measure and momentum dependence, we can now use the powers of the
numerator and denominator factors to label individual integrals.

\subsection{A common notation for an integral family}
We add auxiliary expressions to the physical denominators so that we can
write every required loop-momentum scalar product in terms of the list
$D_1,\ldots,D_N$ and the fixed parameters $\boldsymbol x$. Here $N$ is
the number of expressions in the list, and the auxiliary expressions
allow us to represent numerator factors. Writing the numerator in terms
of this list expresses $\mathcal J$ as a linear combination of integrals
of the form
\begin{equation}
 I_{\boldsymbol a}(\boldsymbol x,\epsilon)
 =\int_{\ell}\prod_{i=1}^{N}D_i^{-a_i},
 \qquad
 \int_{\ell}\equiv
 \prod_{k=1}^{L}\int\frac{\mathrm{d}^{d}\ell_k}{\mathrm{i}\pi^{d/2}},
 \qquad \boldsymbol a=(a_1,\ldots,a_N)\in\mathbb Z^N.
 \label{eq:family}
\end{equation}
Here $\int_\ell$ is shorthand for the loop integration measure and
$\boldsymbol a$ is the vector of $N$ integer powers, with $\mathbb Z$
denoting the set of integers. These powers can be positive, zero or
negative, so we can
write numerator and denominator factors in the equivalent form
\begin{equation}
 I_{\boldsymbol a}(\boldsymbol x,\epsilon)
 =\int_\ell
 \frac{\prod_{i:a_i<0}D_i^{-a_i}}
      {\prod_{i:a_i>0}D_i^{a_i}}.
 \label{eq:index-signs}
\end{equation}
A positive index places $D_i^{a_i}$ in the denominator, a negative index
places $D_i^{-a_i}$ in the numerator, and a zero index removes the factor,
with an empty product equal to one. Auxiliary expressions used only as
numerators therefore have nonpositive indices. An irreducible scalar
product is a numerator scalar product that cannot be written as a linear
combination of the physical denominators and terms depending only on
$\boldsymbol x$~\cite{Weinzierl2022}.

An integral family contains the integrals obtained by varying
$\boldsymbol a$ for a fixed list $D_1,\ldots,D_N$ and fixed kinematics.
Within this family, the target list specifies the integrals that the
amplitude requires, and grouping family members by their positive indices
organizes them into sectors.

\subsection{Sectors and the organization of related integrals}
A sector is a set of integrals with the same positive-index positions,
so they contain the same denominators. We assign a label to each sector
using
\begin{equation}
 \theta_i(\boldsymbol a)=
 \begin{cases}1,&a_i>0,\\0,&a_i\leq0,\end{cases}
 \qquad
 S(\boldsymbol a)=\sum_{i=1}^{N}2^{i-1}\theta_i(\boldsymbol a).
 \label{eq:sector}
\end{equation}
The indicator $\theta_i$ is one when denominator $i$ is present and zero
otherwise, and $S$ is an integer that identifies this set of denominators.
Integrals with different powers of the same denominators belong to the
same sector because $S$ depends only on which indices are positive.
Removing one or more denominators gives a subsector containing a proper
subset of the original sector's denominators.

Scaleless sectors contain no nonzero mass or kinematic scale and therefore
vanish in dimensional regularization, while changes of loop integration
variables can yield symmetry relations between integrals. Kira uses these
zero sectors and symmetry mappings to reduce the number of
unknowns~\cite{KiraOriginal}. We can then use the sector and index labels
to write the identities that relate the remaining family members.

\section{From integration to algebraic reduction}
\label{sec:ibp-reduction}
To express the targets in terms of master integrals, we generate
relations between family members through integration by parts and solve
the resulting equations. Evaluating the master integrals then gives
the target integral values needed to calculate the amplitude.

\subsection{How integration by parts produces relations}
For differentiable functions $f(z)$ and $g(z)$ of a real variable $z$, the
product rule gives the integration identity between endpoints $z_0,z_1$:
\begin{equation}
 \int_{z_0}^{z_1} f'(z)g(z)\,\mathrm{d}z
 =[f(z)g(z)]_{z_0}^{z_1}
 -\int_{z_0}^{z_1} f(z)g'(z)\,\mathrm{d}z.
 \label{eq:ordinary-ibp}
\end{equation}
Primes denote derivatives with respect to $z$, and the brackets denote the
endpoint difference $f(z_1)g(z_1)-f(z_0)g(z_0)$. To apply this product rule
to loop-momentum derivatives, we derive the integration-by-parts (IBP)
identity where the regulated integral has a vanishing boundary term and then
extend it by analytic continuation~\cite{ChetyrkinTkachov,Weinzierl2022}.
For a loop momentum $\ell_k$ and a vector
$v\in\{\ell_1,\ldots,\ell_L,p_1,\ldots,p_E\}$, the identity is
\begin{equation}
 0=\int_{\ell}\frac{\partial}{\partial\ell_k^\mu}
       \left(v^\mu\prod_{i=1}^{N}D_i^{-a_i}\right).
 \label{eq:ibp}
\end{equation}
The index $\mu$ is a spacetime index, and a repeated Lorentz index means
that we sum over components using the spacetime metric. The derivative
acts on both $v^\mu$ and the factors $D_i^{-a_i}$, so the product rule gives
\begin{equation}
 0=\int_{\ell}\left[
 \frac{\partial v^\mu}{\partial\ell_k^\mu}
 -\sum_{i=1}^{N}a_i\,
   \frac{v^\mu(\partial D_i/\partial\ell_k^\mu)}{D_i}
 \right]\prod_{j=1}^{N}D_j^{-a_j}.
 \label{eq:ibp-expanded}
\end{equation}
The first term equals $d$ for $v=\ell_k$ and vanishes when $v$ is an
external momentum or another independent loop momentum. Differentiating
an inverse propagator produces scalar products, which we express through
the family basis to convert each term into an integral with shifted indices:
\begin{equation}
 \sum_{j=1}^{n_\lambda}
 c_{\lambda j}(\boldsymbol a;\boldsymbol x,\epsilon)
 I_{\boldsymbol a+\boldsymbol\Delta_{\lambda j}}
 (\boldsymbol x,\epsilon)=0.
 \label{eq:ibp-linear}
\end{equation}
Here $\lambda$ identifies the chosen pair $(k,v)$, $n_\lambda$ is the number of
terms, and $\boldsymbol\Delta_{\lambda j}$ is the integer change in the
index vector for term $j$. Differentiation gives the coefficient
$c_{\lambda j}$, which depends on the integral indices and the parameters
$\boldsymbol x,\epsilon$~\cite{ChetyrkinTkachov}.

The $L$ choices of differentiated momentum and $L+E$ choices of $v$ give
$L(L+E)$ basic identities at each index vector, some of which may depend
linearly on the others. Further relations follow from Lorentz invariance
because a common rotation or Lorentz boost of all external momenta
preserves a scalar integral, giving the infinitesimal identity
\begin{equation}
 \sum_{\beta=1}^{E}\left(
 p_\beta^\nu\frac{\partial}{\partial p_{\beta\mu}}
 -p_\beta^\mu\frac{\partial}{\partial p_{\beta\nu}}
 \right) I_{\boldsymbol a}=0.
 \label{eq:li}
\end{equation}
Here $\mu,\nu$ are free spacetime indices, meaning that the identity holds
for each pair of components, and the derivatives act on the
external-momentum dependence of the integral. Contracting these indices
with external momenta gives scalar relations~\cite{GehrmannRemiddi}
that join the IBP and symmetry identities to form the reduction system.

\subsection{From relations to master integrals}
Let $n_I$ be the number of distinct integrals in the generated equations
and $n_{\mathrm{eq}}$ the number of equations. We arrange these integrals
in a column vector so that the system takes the matrix form
\begin{equation}
 A(\boldsymbol x,\epsilon)\,\boldsymbol I=0,
 \qquad
 \boldsymbol I=(I_{\boldsymbol a^{(1)}},\ldots,
 I_{\boldsymbol a^{(n_I)}})^{\mathsf T}.
 \label{eq:matrix}
\end{equation}
The superscript $\mathsf T$ denotes transpose, so $\boldsymbol I$ is a
column vector of integrals whose dependence on $\boldsymbol x,\epsilon$
we leave implicit here. The matrix $A$ has $n_{\mathrm{eq}}$ rows and
$n_I$ columns, with each entry $A_{ij}$ giving the coefficient of integral
$j$ in equation $i$. Solving this system expresses dependent integrals in
terms of an independent set. This set is a master basis if the reduction
relations express every required integral as a linear combination of its
members. For the standard class of Feynman integrals at fixed graph,
this basis contains a finite number
of masters~\cite{SmirnovPetukhov}. We write them as
$M_1,\ldots,M_{n_M}$, where $n_M$ is the number of master integrals, so
the reduction formula is
\begin{equation}
 I_{\boldsymbol a}(\boldsymbol x,\epsilon)
 =\sum_{m=1}^{n_M}C_{\boldsymbol a m}(\boldsymbol x,\epsilon)
 M_m(\boldsymbol x,\epsilon).
 \label{eq:reduction}
\end{equation}
Here $C_{\boldsymbol a m}$ is the coefficient of master $M_m$ in the
reduction of target $I_{\boldsymbol a}$. The solver obtains it by adding,
multiplying and dividing the identity coefficients, so it is a rational
function, meaning a ratio of polynomials in $\boldsymbol x,\epsilon$. Different
independent sets can define different master bases, and an insufficient
equation set can leave additional integrals unresolved. Establishing the
master basis therefore requires checking that the generated equations
are enough to reduce the required integrals.

Equation~\eqref{eq:reduction} gives the coefficients multiplying the
masters, whose evaluation requires a further calculation, for example by
solving differential equations with known values or limits as boundary
conditions~\cite{GehrmannRemiddi}. Kira provides the reduction rules
needed for this calculation by generating and processing the identities
through the procedure described next.

\section{How Kira performs the reduction}
\label{sec:kira-procedure}
Kira implements the Laporta approach by generating a finite set of
identities and eliminating integrals according to an
ordering~\cite{Laporta}. To obtain the coefficients in
Eq.~\eqref{eq:reduction}, Kira selects seeds, removes redundant
equations and solves the retained system. We examine how Kira performs
these steps to specify the seed and elimination choices that the
proposed models would make.

\subsection{Kira's input and seed generation}
The user supplies the families in Eq.~\eqref{eq:family}, relations among
the parameters $\boldsymbol x$, the target list, and bounds on equation
generation so that Kira can initialize the families and sectors. During
this initialization, Kira detects scaleless sectors and identifies symmetry
relations and mappings between compatible families~\cite{KiraOriginal}.

A seed specifies the index vector $\boldsymbol a$ at which Kira evaluates
Eq.~\eqref{eq:ibp-linear} to generate several derivative identities, and
the seed set may include integrals beyond the target list. The shifts
$\boldsymbol\Delta_{\lambda j}$ in these identities can introduce further
integrals outside the seed set, so the resulting system depends on both
the chosen seeds and the terms generated from them. To set limits on the
indices of the seeds, Kira uses the quantities
\begin{align}
 t(\boldsymbol a)&=\sum_{i=1}^{N}\theta_i(\boldsymbol a),
 &r(\boldsymbol a)&=\sum_{a_i>0}a_i,\\
 s(\boldsymbol a)&=-\sum_{a_i<0}a_i,
 &n_{\mathrm{dot}}(\boldsymbol a)&=r(\boldsymbol a)-t(\boldsymbol a).
 \label{eq:complexities}
\end{align}
Here $t$ is the number of distinct denominators, $r$ is the sum of their
powers, and $s$ is the sum of the numerator powers. The quantity
$n_{\mathrm{dot}}$ is the sum of the extra denominator powers beyond
one for each denominator. We use this symbol to keep $d$ reserved for the
integration dimension, and we restrict equation generation by specifying
the allowed sectors and limits
\begin{equation}
 r\leq r_{\max},\qquad s\leq s_{\max},\qquad
 n_{\mathrm{dot}}\leq n_{\mathrm{dot},\max},
 \label{eq:seed-bounds}
\end{equation}
where $r_{\max}$, $s_{\max}$ and $n_{\mathrm{dot},\max}$ are the upper
limits on the corresponding quantities. The user sets $r_{\max}$ and
$s_{\max}$ and may also set $n_{\mathrm{dot},\max}$. Because these choices determine which seeds
contribute equations, a reduction run must establish whether the bounds
supply enough relations for the targets~\cite{Kira12}.

Kira 3 changes seeding in symmetry-related sectors and provides the optional
\texttt{truncate\_sp} setting to restrict numerator seeding in subsectors,
with parameter choices that require checks against the expected master
basis~\cite{Kira3}. The generated system can still contain redundant
equations, which elimination identifies as it processes the relations
into reduction rules.

\subsection{Integral ordering, sector order and forward elimination}
Elimination solves an identity for one integral and substitutes the result
into other identities, using an integral ordering to select which unknown
to eliminate from each equation. Kira defines these orderings through
sector labels and the powers of the numerator and denominator factors,
with choices that prioritize numerator powers or repeated denominators
and thereby affect the master basis and
runtime~\cite{Kira12}.

To describe this procedure, write $I_j=I_{\boldsymbol a^{(j)}}$ for the
integrals in Eq.~\eqref{eq:matrix} and choose the ordering
$I_1\succ I_2\succ\cdots\succ I_{n_I}$, where $\succ$ gives priority
for elimination. In a row whose first nonzero term contains $I_k$, let
$b_j$ denote the coefficient of $I_j$, so choosing the nonzero coefficient
$b_k$ as the pivot gives
\begin{equation}
 b_k I_k+\sum_{j>k}b_j I_j=0
 \quad\Longrightarrow\quad
 I_k=-\sum_{j>k}\frac{b_j}{b_k}I_j.
 \label{eq:pivot}
\end{equation}
The coefficients $b_j$ depend on $\boldsymbol x,\epsilon$, and substituting
this expression into the other rows removes their dependence on $I_k$.
Repeating this step gives forward elimination, which builds a triangular
system in which each pivot integral depends on integrals later in the
ordering. Back substitution then replaces the remaining pivot integrals
using their own rules, completing the Laporta solution~\cite{Laporta}.

Sector scheduling sets the order in which the solver processes groups of
integrals or equations, taking into account which other integrals their
equations contain. Ordinary IBP identities for a sector involve that
sector and its subsectors because nonpositive indices remain nonpositive
when applying the derivative. Equations for one sector can therefore
depend on integrals from its subsectors. The Kira 3 description starts
equation generation with sectors earlier in its ordering and inserts
previously obtained pivot relations into new equations~\cite{Kira3},
with the order of processing depending on the stage of the reduction.

When several rows share the same first integral, the choice of pivot row
changes the substitutions that follow. Choosing row $\rho$ to eliminate
column $k$, for example, updates row $i$ by
\begin{equation}
 A_{ij}\ \longleftarrow\ A_{ij}
 -\frac{A_{ik}}{A_{\rho k}}A_{\rho j},
 \qquad A_{\rho k}\ne0.
 \label{eq:row-update}
\end{equation}
Here $i$ is the index of a row other than the pivot row, so $i\ne\rho$,
and $j$ is the column index.
Subtracting the stated multiple of row $\rho$ sets the coefficient in
column $k$ to zero and can create nonzero entries in other columns,
a change called fill-in when a previously zero entry becomes nonzero.
These substitutions also make the coefficient expressions longer by
combining their numerators and denominators, so elimination can require
more time and memory as it proceeds. Redundant equations add to this
work~\cite{Kant}, so Kira uses modular elimination to
select equations before processing their full symbolic coefficients.

\subsection{Removing redundant equations and selecting target dependencies}
To reduce the number of equations it must solve, Kira identifies rows
that follow from other equations. Elimination reduces every coefficient
in such a row to zero, allowing the solver to omit that row. Kira
tests these dependencies with modular arithmetic~\cite{KiraOriginal}
by substituting sample values $\boldsymbol x_0,\epsilon_0$ and mapping
the coefficients to a finite field $\mathbb F_P$ for a prime integer $P$:
\begin{equation}
 A(\boldsymbol x,\epsilon)
 \longmapsto A(\boldsymbol x_0,\epsilon_0)\bmod P.
 \label{eq:modular-selection}
\end{equation}
Integers from $0$ to $P-1$ represent the elements of $\mathbb F_P$, and
addition and multiplication keep results in this range by taking the
remainder upon division by $P$. Division uses the modular inverse of a
nonzero element, which multiplies that element to one modulo $P$, so the
sampling procedure excludes points where a required denominator vanishes.
This arithmetic keeps field elements bounded in size while allowing the
solver to identify independent rows and retain their original symbolic
identities.

The matrix rank is the largest number of independent rows, where no row
is a linear combination of the others in the chosen set. Some sample
values can make otherwise independent rows dependent, giving a rank
smaller than that of the symbolic matrix. Checks must account for this
possibility when using the probabilistic equation-selection method
developed by Kant~\cite{Kant}.

Kira also tracks which independent equations contribute to the target
reductions, using the dependencies of the forward-eliminated system for
this selection in Kira 3~\cite{Kira3}. Seed selection therefore supplies
the integrals from which Kira generates equations, and equation selection
then retains the relations needed for the targets so that the solver can
perform the full reduction on this smaller system.

\subsection{Back substitution and the final symbolic result}
To obtain $C_{\boldsymbol a m}$ from the retained equations, Kira applies
back substitution to the pivot rules produced by forward elimination.
Since these rules can contain other pivot integrals, back substitution
inserts their expressions recursively until the right-hand sides contain
only masters, giving Eq.~\eqref{eq:reduction}. Kira's original symbolic
solver carries out this calculation using rational-function arithmetic
through the computer algebra system Fermat~\cite{KiraOriginal}.

Kira also solves systems at finite-field sample points and reconstructs the
rational coefficients with FireFly, the reconstruction library introduced
into the solver in Kira 2.0~\cite{Kira2}. In this procedure, each sampled
coefficient has the form
\begin{equation}
 C_{\boldsymbol a m}(\boldsymbol x_0,\epsilon_0)\bmod P.
 \label{eq:sampled-coefficient}
\end{equation}
The modular operations use exact arithmetic, and reconstruction recovers
$C_{\boldsymbol a m}(\boldsymbol x,\epsilon)$ from sufficient samples
before checking the result at additional points. Reconstruction therefore
forms part of the calculation of the symbolic reduction coefficients,
after which the master integrals still require evaluation~\cite{Kira2}.

Because seed selection determines the equations that enter this procedure
and elimination decisions determine the substitutions and intermediate
expressions, the proposed ML models would guide these two operations to
reduce the time and memory needed for the reduction.


\section{ML for seed selection and elimination ordering}
\label{sec:ml-roles}
To reduce the time and memory needed for IBP reduction, we will study
two System 1 models: one for seed selection and one for elimination
ordering. Here a System 1 model is a small machine learning model that
learns from recorded reductions and makes a choice from the information
available at a given stage. We would add the seed model before equation
generation and the elimination model at the stage where the solver
selects a pivot. Each model would assign scores to the available choices,
and Kira would perform the algebra for the selected choice. Training
these models would require target integrals with known master expansions,
the seeds that supply their reduction equations, and the steps used to
solve those equations.

\subsection{Selecting seeds before equation generation}
To select seeds, the first model would receive the family definition,
the target index vectors $\boldsymbol a$, the chosen master basis and
a set of candidate seeds within specified sectors and index bounds.
For each candidate, it would use the seed indices and the target
information to estimate how useful the resulting equations would be
for the requested reduction.

During training, we would pair each target with the seed sets recorded
in its known reductions. These examples would teach the model to rank
seeds that together provide enough equations to reduce the target.
Different reductions can use different sufficient seed sets, so the
training data would retain these alternatives and their equation counts.
When applying the trained model, we would select highly ranked seeds
and ask Kira to generate their identities. If the equations leave a
target unresolved in the chosen master basis, we would add further
seeds and repeat the reduction. The second model would guide the
elimination of the resulting equations.

\subsection{Choosing elimination steps within the solver}
To limit fill-in and the growth of coefficient expressions, the second
model would rank the allowed pivot choices using the current equations,
their nonzero coefficients and the integral ordering. For each recorded
reduction step, the training example would contain the equations before
the choice, the available pivots, the selected row and integral, and the
operations that followed. The remaining recorded steps and their
operation counts would show how that choice contributed to a complete
reduction.

During a reduction, the model would select an allowed pivot and Kira
would carry out the row operations in Eq.~\eqref{eq:row-update}. The
model would then receive the updated system and select the next pivot,
continuing until Kira can complete the reduction by back substitution.
Each choice must respect the integral ordering and chosen master basis.
Preparing the model inputs and obtaining its scores must take less time
than its choices save in the algebraic calculation. Kira would use its
existing reconstruction procedure for calculations over finite fields.

\subsection{Integral data required for training and testing}
To train and test these models, we would require examples that specify
the family denominators, kinematic parameters, target integrals and
master basis, together with the exact coefficients in
Eq.~\eqref{eq:reduction}. Each example would also contain its source
IBP equations, their seeds and derivative operators, and any symmetry
relations used. These records would connect the seed model's choices
to the equations available to the elimination model.

We would record both the operations used to construct each target and
the elimination and back-substitution steps used to reduce it. The
records would include intermediate equations, pivot choices and
operation counts, so the elimination model can learn from complete
sequences. The stored master coefficients would supply the reference
answer for checking that each target reduces completely in the chosen
basis, a requirement discussed in Section~\ref{sec:ibp-reduction} and
the Kira documentation~\cite{Kira3}. During tests, the models would
receive the problem information available at their respective stages;
the recorded solutions would remain separate for checking the results.
Section~\ref{sec:integrals-generation} explains how we would construct
the target integrals and obtain these records together.

\section{Target Integral Generation for Machine Learning}
\label{sec:integrals-generation}
To train and test the two models described in Section~\ref{sec:ml-roles},
we will construct target integrals from master integrals and retain the
identities, seeds and algebraic steps used in each construction. We have
completed the first stage by identifying master bases for 16 one-loop
and two-loop families with Kira 3.1. These bases contain 49 entries when
counted separately for each family. The following subsections define
the families and explain how we obtained and checked their bases,
before describing the planned construction of targets from these
masters. The completed calculation determines symbolic integral
representatives; target construction and the calculation of their
symbolic reduction coefficients remain subsequent stages.

\subsection{One-loop families and their kinematics}
We use the integral definition and normalization in
Eq.~\eqref{eq:family}, with $d=4-2\epsilon$ and inverse propagators
$D_i=q_i^2-m_i^2+\mathrm{i}0$. For each family, we specify the internal
momenta $q_i$, the masses $m_i$ and the scalar products of the external
momenta, so every index vector defines a particular integral. In the
equations below we suppress the common $\mathrm{i}0$ prescription.
We write $\mu=m^2$ for a common nonzero squared mass and
$\mu_1=m_1^2$, $\mu_2=m_2^2$ for two independent nonzero squared masses.
The symbols $x_s$ and $x_t$ denote momentum invariants and remain
distinct from the index sums $s$ and $t$ in
Eq.~\eqref{eq:complexities}.

The one-loop collection begins with a massive tadpole, whose single
inverse propagator is $D_1=\ell_1^2-\mu$ and which has no external
momentum. We then include four bubble families with external momentum
$p$ and inverse propagators
\begin{equation}
 D_1=\ell_1^2-m_1^2,\qquad
 D_2=(\ell_1-p)^2-m_2^2,\qquad p^2=x_s.
 \label{eq:generated-bubble-family}
\end{equation}
The squared-mass pairs are $(0,0)$, $(\mu,0)$, $(\mu,\mu)$ and
$(\mu_1,\mu_2)$, corresponding to a massless bubble, a bubble with one
massive line, an equal-mass bubble and an unequal-mass bubble.
We keep $x_s$ as an independent nonzero scale, with the massive
families defined at generic values of their independent parameters.

To include three-point integrals, we define triangle families by
\begin{equation}
 \begin{aligned}
 D_1&=\ell_1^2-m_1^2,&
 D_2&=(\ell_1-p_1)^2-m_2^2,&
 D_3&=(\ell_1-p_1-p_2)^2-m_3^2,\\
 p_1^2&=u_1,&p_2^2&=u_2,&(p_1+p_2)^2&=u_3.
 \end{aligned}
 \label{eq:generated-triangle-family}
\end{equation}
Here $u_1,u_2,u_3$ are the squared external momenta, also called
external virtualities, and
$p_1\cdot p_2=(u_3-u_1-u_2)/2$ fixes the remaining external scalar
product. We use three massless families with
$(u_1,u_2,u_3)=(0,0,x_s)$, $(0,x_2,x_3)$ and $(x_1,x_2,x_3)$.
The parameters $x_1,x_2,x_3$ denote independent nonzero virtualities.
We also use two families with $m_1^2=m_2^2=m_3^2=\mu$ and
virtualities $(0,0,x_s)$ and $(x_1,x_2,x_3)$. Thus the five triangle
families vary both the internal masses and the number of nonzero
external virtualities.

For four-point integrals, we use three massless box families with
$p_4=-p_1-p_2-p_3$ and
\begin{equation}
 \begin{aligned}
 D_1&=\ell_1^2,&D_2&=(\ell_1-p_1)^2,\\
 D_3&=(\ell_1-p_1-p_2)^2,&
 D_4&=(\ell_1-p_1-p_2-p_3)^2,\\
 x_s&=(p_1+p_2)^2,&x_t&=(p_2+p_3)^2.
 \end{aligned}
 \label{eq:generated-box-family}
\end{equation}
Writing $u_i=p_i^2$ for $i=1,\ldots,4$, we fix the independent
external scalar products through
\begin{equation}
 p_1\cdot p_2=\frac{x_s-u_1-u_2}{2},\qquad
 p_2\cdot p_3=\frac{x_t-u_2-u_3}{2},\qquad
 p_1\cdot p_3=\frac{u_4+u_2-x_s-x_t}{2}.
 \label{eq:generated-box-kinematics}
\end{equation}
The virtuality choices are $(0,0,0,0)$, $(0,0,0,x_4)$ and
$(x_1,x_2,x_3,x_4)$, where $x_4$ is the fourth nonzero virtuality
parameter. These choices give boxes with zero, one and four off-shell
legs; in this classification an on-shell external leg has $p_i^2=0$.
We treat the remaining invariants as independent subject to the stated
relations. These definitions give 13 one-loop families, which we
extend with the following three two-loop families.

\subsection{Two-loop families and auxiliary expressions}
The two-loop families use loop momenta $\ell_1,\ell_2$. For a massless
two-point integral with external momentum $p$ and $p^2=x_s$, we take
\begin{equation}
 (D_1,D_2,D_3,D_4,D_5)
 =\bigl(\ell_1^2,\ell_2^2,(\ell_1-\ell_2)^2,
        (\ell_1-p)^2,(\ell_2-p)^2\bigr).
 \label{eq:generated-two-loop-propagator}
\end{equation}
These five expressions determine the five loop-dependent scalar
products $\ell_1^2,\ell_2^2,\ell_1\cdot\ell_2,\ell_1\cdot p$ and
$\ell_2\cdot p$. Varying their indices therefore includes numerator
insertions as well as repeated propagators.

For the equal-mass sunrise with the same external invariant, we use
\begin{equation}
 \begin{aligned}
 D_1&=\ell_1^2-\mu,&D_2&=\ell_2^2-\mu,&
 D_3&=(p-\ell_1-\ell_2)^2-\mu,\\
 D_4&=(p-\ell_1)^2,&D_5&=(p-\ell_2)^2,&
 a_4,a_5&\leq0.
 \end{aligned}
 \label{eq:generated-sunrise-family}
\end{equation}
The first three expressions are the physical inverse propagators.
The auxiliary expressions $D_4,D_5$ complete the five scalar products
needed to express the derivatives in the IBP identities. Their
nonpositive indices restrict them to numerator factors or their
absence, so the top sector contains only $D_1,D_2,D_3$ as
denominators. We keep $x_s$ and $\mu$ independent.

The final family is the equal-mass two-loop vacuum sunset, which has
no external momentum and uses
\begin{equation}
 D_1=\ell_1^2-\mu,\qquad D_2=\ell_2^2-\mu,\qquad
 D_3=(\ell_1-\ell_2)^2-\mu.
 \label{eq:generated-vacuum-family}
\end{equation}
These three expressions determine all loop scalar products for this
vacuum family. For every family, we include the sector containing all
physical propagators and all its subsectors, so the resulting basis
also accounts for integrals with removed propagators. The number of
physical propagators in this top sector is denoted by $t_{\rm top}$.
It is one, two, three and four for the tadpole, bubble, triangle and
box families, respectively, five for the massless two-loop propagator,
and three for the sunrise and vacuum families. This number sets the
initial range of identities used to identify the masters.

\subsection{Identification and checks of the master bases}
For each family, we used Kira 3.1 to generate IBP relations together
with the applicable Lorentz-invariance and symmetry relations.
Kira removes scaleless sectors and uses finite-field elimination
to identify the integrals that remain independent in the equations
needed for a selected integral range~\cite{KiraOriginal,Kira3}.
We used Kira's integral ordering~1 throughout, which orders sectors
by their labels and prefers numerator insertions over extra
denominator powers when selecting representatives within a
sector~\cite{Kira12}. Keeping this ordering fixed gives a specified
basis against which to compare later calculations.

We first generated identities from seeds satisfying
\begin{equation}
 r\leq t_{\rm top}+2,\qquad s\leq2,\qquad
 n_{\rm dot}\leq2,
 \label{eq:master-discovery-seeds}
\end{equation}
and selected integrals to probe the basis within
\begin{equation}
 r\leq t_{\rm top}+1,\qquad s\leq1,\qquad
 n_{\rm dot}\leq1.
 \label{eq:master-discovery-probes}
\end{equation}
Here $r$ is the sum of positive indices, $s$ is minus the sum of
negative indices and $n_{\rm dot}=r-t$ is the number of extra
denominator powers, as defined in Eq.~\eqref{eq:complexities}.
Both sets of limits apply to the top sector and its subsectors.
The probe integrals specify the range for which Kira must identify
the required masters, while the wider seed range supplies relations
that can eliminate the additional integrals introduced by index
shifts. These probes serve the basis calculation; the ML examples
will come from the subsequent reverse construction.

Kira evaluated the coefficients of these relations over a finite
field using the substitution in Eq.~\eqref{eq:modular-selection}.
For the first calculation we used the prime $P_1=2147483647$ and
dimension value $d_1=37$. To specify the kinematic substitutions,
define the ordered parameter list
\begin{equation}
 \boldsymbol z=(x_s,x_t,x_1,x_2,x_3,x_4,\mu,\mu_1,\mu_2).
 \label{eq:master-sample-parameters}
\end{equation}
The first sample was
\begin{equation}
 \boldsymbol z^{(1)}=(-13,-19,-11,-17,-29,-43,7,7,23),
 \label{eq:master-first-sample}
\end{equation}
using only the entries present in each family and retaining all
zeros and equal-mass conditions in its definition. We interpret
these assignments modulo $P_1$ when evaluating the identity
coefficients. The dimension sample corresponds to
$\epsilon_1=(4-d_1)/2$ in the same field. These substitutions
test algebraic relations; the master labels continue to refer to
integrals at symbolic $\boldsymbol x,\epsilon$.

We then enlarged the seed and probe ranges to check whether the
initial basis remained sufficient:
\begin{equation}
 \begin{aligned}
 \text{seeds:}\quad&
 r\leq t_{\rm top}+3,\qquad s\leq3,\qquad n_{\rm dot}\leq3,\\
 \text{probes:}\quad&
 r\leq t_{\rm top}+2,\qquad s\leq2,\qquad n_{\rm dot}\leq2.
 \end{aligned}
 \label{eq:master-expanded-bounds}
\end{equation}
We repeated the calculation first at the original sample and then
at a second sample with prime $P_2=1000000007$, dimension $d_2=41$
and
\begin{equation}
 \boldsymbol z^{(2)}=(-31,-47,-23,-37,-53,-71,11,11,29).
 \label{eq:master-second-sample}
\end{equation}
Changing the prime, the dimension and the independent kinematic
parameters provides a second check of the linear dependencies,
since particular substitutions can lower the rank of the
symbolic coefficient matrix~\cite{Kant}.

In both enlarged calculations, we supplied the original masters as
the preferred basis and enabled Kira's check for additional or
unused masters~\cite{Kira3}. We required the complete list of
master index vectors to agree in all three calculations, and
also checked each vector against its family definition and sector
label. All 16 families satisfied these requirements, giving
48 completed master-identification calculations. The agreement
provides evidence for basis stability over the tested index ranges
and sample points. Completeness for arbitrary integer indices
would require further relations or an independent argument, and
special kinematic limits can change the basis.

We stopped each calculation after master identification. This stage
determines which symbolic integrals form the starting basis within
the checked domains; obtaining the rational reduction coefficients
for new targets will require the algebraic construction described
below. The numerical values and $\epsilon$ expansions of the
master integrals remain separate from this basis determination.

\subsection{Resulting bases and the starting coefficients}
The massive tadpole gave one master. The massless, one-mass,
equal-mass and unequal-mass bubbles gave one, two, two and three
masters, respectively. The massive bubble bases include the
necessary tadpole subsectors, with two distinct tadpoles when
the internal masses differ.

The massless triangles with one, two and three nonzero external
virtualities gave one, two and four masters, respectively.
For the first two families, Kira chose bubble integrals in
subsectors as the masters. The three-virtuality family also
retained a triangle integral. The equal-mass triangles with
one and three nonzero virtualities gave three and five masters.
The massless boxes with zero, one and four off-shell legs gave
three, four and eleven masters, respectively, including the
required lower sectors in each count.

The massless two-loop propagator gave two masters, the equal-mass
sunrise gave three, and the equal-mass vacuum sunset gave two.
For the denominator order in
Eq.~\eqref{eq:generated-sunrise-family}, the sunrise basis is
\begin{equation}
 M_1=I_{(1,1,0,0,0)},\qquad
 M_2=I_{(1,1,1,0,0)},\qquad
 M_3=I_{(1,1,1,-2,0)}.
 \label{eq:generated-sunrise-basis}
\end{equation}
The first integral factorizes into a product of two massive
tadpoles because removing $D_3$ separates the loop integrations.
The other two integrals retain all three sunrise propagators,
and $M_3$ contains the numerator factor $D_4^2$. This example
shows how the chosen ordering can retain a numerator integral
as a master.

The counts sum to 49 master entries across the 16 families.
We count each basis within its own family, so shared tadpole
and bubble subsectors can appear in more than one family.
The complete family definitions and master index vectors are
available with the accompanying data~\cite{MasterIntegralData}.
Each entry specifies the denominators, kinematics, index vector,
sector and mass dimension. With the normalization in
Eq.~\eqref{eq:family}, the latter is
\begin{equation}
 [I_{\boldsymbol a}]_{\rm mass}
 =L d-2\sum_{i=1}^{N}a_i,
 \label{eq:generated-mass-dimension}
\end{equation}
where $[I_{\boldsymbol a}]_{\rm mass}$ is the power of mass
carried by the integral. Each loop measure contributes $d$
powers and each factor $D_i^{-a_i}$ contributes $-2a_i$.
This dimension will constrain the coefficients used to combine
integrals in the next stage.

For that construction, we will keep the identified masters
symbolic and initialize each master as itself. If
$\boldsymbol\mu^{(m)}$ is the index vector of master $M_m$,
the starting coefficients are
\begin{equation}
 M_m=I_{\boldsymbol\mu^{(m)}},
 \qquad
 C_{\boldsymbol\mu^{(m)}n}=\delta_{mn},
 \qquad m,n=1,\ldots,n_M.
 \label{eq:master-initialization}
\end{equation}
Here $n_M$ is the number of masters in the selected family,
and $\delta_{mn}$ is one when $m=n$ and zero otherwise.
Thus each master has coefficient one for itself and zero
for the others. We will hold the family definition and basis
fixed while constructing its targets, starting with subsets
of these masters and extending the collection through the
identities below. This gives each subsequent target a specified
reference basis and an exact master expansion.

\subsection{Linear combinations with known master expansions}
To construct target expressions with known master expansions, we would
combine family integrals whose reductions we have already obtained.
The initial choices would be the masters themselves, and later choices
would also include integrals generated through the identities described
below.
Let $I_{\boldsymbol a^{(1)}},\ldots,I_{\boldsymbol a^{(n_T)}}$ be the
selected integrals, where $n_T$ is their number. Choosing weights
$w_j(\boldsymbol x,\epsilon)$ gives a target expression $\mathcal T$
whose master expansion follows by substitution. Suppressing the common
arguments $\boldsymbol x,\epsilon$, we obtain
\begin{equation}
 \mathcal T
 =\sum_{j=1}^{n_T}w_j I_{\boldsymbol a^{(j)}}
 =\sum_{m=1}^{n_M}C^{\mathcal T}_m M_m,
 \qquad
 C^{\mathcal T}_m
 =\sum_{j=1}^{n_T}w_j C_{\boldsymbol a^{(j)}m}.
 \label{eq:generated-combination}
\end{equation}
Here $j$ identifies a selected integral, $m$ identifies a master, and
$C^{\mathcal T}_m$ is the coefficient of that master in $\mathcal T$.
The coefficients $C_{\boldsymbol a^{(j)}m}$ are the known reductions
defined in Eq.~\eqref{eq:reduction}, so multiplication and addition give
the exact answer for the new expression.

The weights may be constants or rational functions of the masses,
external momentum invariants and $\epsilon$. They remain independent of
the loop momenta, and we would choose their mass dimensions so that the
terms in the sum have the same mass dimension. This dependence is
consistent with the coefficient functions in IBP
reduction~\cite{Kira3}. Varying the selected integrals and their weights
would produce different target expressions, including expressions with
cancellations between master coefficients. The corresponding seed and
elimination records would include the reductions of the selected
integrals and the operations used to combine their results.

Equation~\eqref{eq:generated-combination} defines a linear combination
of family integrals. Identifying that expression with one integral
$I_{\boldsymbol b}$ requires an identity involving the specific index
vector $\boldsymbol b$. We would therefore use IBP relations to
determine combinations that correspond to individual family members,
as described next. This step supplies targets in their integral form
and a separate record of their known master expansions.

\subsection{Constructing individual targets through IBP identities}
To construct a new individual target with a known reduction, we would
use IBP identities to express that integral in terms of integrals whose
reductions are already known. We would generate the identities in
Eq.~\eqref{eq:ibp-linear} from a seed vector $\boldsymbol\sigma$ and
record the differentiated loop momentum and the vector $v$ used in
Eq.~\eqref{eq:ibp}. The pair $(k,v)$ specifies
the IBP operator, so the seed and this pair determine the source of each
equation~\cite{KiraOriginal}. Consider an identity that contains a new
integral $I_{\boldsymbol b}$ and $n_R$ integrals whose reductions we
already know:
\begin{equation}
 h_0 I_{\boldsymbol b}
 +\sum_{j=1}^{n_R}h_j I_{\boldsymbol a^{(j)}}=0,
 \qquad h_0\ne0.
 \label{eq:generation-identity}
\end{equation}
Here $n_R$ is the number of known integral terms, and $h_0,h_j$ are the
coefficients obtained from the identity, with their dependence on
$\boldsymbol x,\epsilon$ implicit. Solving for $I_{\boldsymbol b}$ and
substituting the known master expansions gives
\begin{equation}
 \begin{aligned}
 I_{\boldsymbol b}
 &=-\sum_{j=1}^{n_R}\frac{h_j}{h_0}
        I_{\boldsymbol a^{(j)}},\\
 C_{\boldsymbol b m}
 &=-\sum_{j=1}^{n_R}\frac{h_j}{h_0}
        C_{\boldsymbol a^{(j)}m}.
 \end{aligned}
 \label{eq:generated-integral-reduction}
\end{equation}
The identity therefore determines both the new integral and its
coefficients in the master basis. We would add this integral to the
collection of known reductions and use it in later constructions,
retaining the conditions under which division by $h_0$ is valid.
Starting with the coefficients in Eq.~\eqref{eq:master-initialization}
and applying Eq.~\eqref{eq:generated-integral-reduction} repeatedly
would build each target from previously known integrals. To recover
its reduction, we would follow these recorded relations from the target
back to the masters and substitute the corresponding expressions.

Some identities contain several integrals whose reductions are still
unknown. In that case, we would combine enough relations to solve for
those integrals before extending the collection, retaining a target
only when these relations determine its complete master expansion.
This may require solving a coupled system before any of its new
integrals can enter the collection. Changing the seeds and
IBP operators would generate targets in different sectors and with
different numerator and denominator powers. We would record the
quantities $t,r,s,n_{\mathrm{dot}}$ from Eq.~\eqref{eq:complexities}
to describe this variation. Each accepted construction must give a
reduction entirely in the chosen master basis, with the exact
coefficients retained as its reference answer.

\subsection{Records for seed selection and elimination}
To train the models on the choices made during reduction, we would store
the source equations and the operations used to reduce each target.
Each record would include the family, the target index vector or target
expression, and its exact master coefficients. For every source IBP
equation, we would store the seed $\boldsymbol\sigma$, the operator
$(k,v)$ and the original equation before substitutions. If a construction
also uses a symmetry relation, we would record that relation and its
integral mapping. These records would identify the seeds and identities
needed to reproduce each equation used in the construction.

The seed record of a new target must include the seeds used for the
integrals that appear in its construction. For the single-identity
step in Eq.~\eqref{eq:generation-identity}, let
$\mathcal S_{\boldsymbol a}$ be the stored seed set for integral
$I_{\boldsymbol a}$. The new record is then
\begin{equation}
 \mathcal S_{\boldsymbol b}
 =\{\boldsymbol\sigma\}\,\cup\!
   \bigcup_{j=1}^{n_R}\mathcal S_{\boldsymbol a^{(j)}},
 \qquad
 \mathcal S_{\boldsymbol\mu^{(m)}}=\varnothing.
 \label{eq:generated-seed-record}
\end{equation}
The union collects the seed of the new identity and the seeds used in
all its earlier reductions, keeping each seed once. The empty set
$\varnothing$ applies to a starting master because expressing it as
itself requires no reduction equations. When a construction uses
several new identities, we would include every seed used to generate
them. We would also retain the individual equations and their
dependencies, since the seed set alone does not specify the reduction
sequence.

We would also store the candidate seeds considered during construction
and the sector and index bounds used to generate them. The collected
seeds would provide a set known to supply enough equations for each
accepted target. If different constructions produce the same target,
we would retain their seed sets together with the corresponding
reduction records. For a target expression $\mathcal T$, we would
combine the seed records of its constituent integrals to obtain an
initial sufficient set, which cancellations in the combined expression
may allow us to reduce further.

For the elimination model, we would replay the source equations in an
allowed forward reduction order and record the system before each
decision. The record would contain the available pivot choices, the
selected equation, the integral eliminated, and the row operations and
substitutions that follow. We would require these steps to respect the
chosen integral ordering and master basis, using the elimination and
back-substitution operations described in
Section~\ref{sec:kira-procedure}~\cite{Laporta}. The target's final
master expansion would remain the reference answer for checking the
complete sequence. This replay would turn the construction record into
examples of the decisions that the model will encounter inside the
solver.

We would record the operations used to construct expressions and the
operations used to reduce the source equations separately. For each
reduction sequence, we would state how we count additions, subtractions,
multiplications and divisions of coefficients, and record changes in
the number of nonzero matrix entries. Each example would therefore
include the reduction sequence, its operation counts and its exact
result. Constructing a target in this way gives a known reduction
without requiring a search through every possible reduction sequence.
We would identify a seed set or sequence as minimal only when a
mathematical argument establishes that it uses the fewest seeds or
operations under the stated rules.

\subsection{Generated data for training and tests}
To test the models on generated targets excluded from training, we would
divide the constructed examples into separate training and test
collections and keep the test solutions hidden from the models.
Targets related by symmetry, integral relabeling or different
construction histories would remain in the same collection, so a
stored version of a test target cannot also become a training example.
We would reserve complete families to measure performance on families
absent from training, and include generated targets with larger powers
and higher loop orders to examine performance on larger reductions.

The seed model would learn from the recorded sufficient seed sets, and
the elimination model would learn from the recorded equations and
reduction steps. During tests, each model would receive only the
corresponding problem information, and the exact master coefficients
would provide the reference for checking its result. For examples with
an established minimum, the recorded operation counts would allow a
comparison with that minimum. The remaining examples would provide
known reductions against which to compare the model's choices.

Before including an example in either collection, we would check its
identities and reproduce its master expansion from the stored source
equations. Keeping the original equations and their construction
history would make these checks possible for both individual integrals
and linear combinations of integrals.

\clearpage
\renewcommand{\refname}{Citations}
\bibliographystyle{unsrt}
\bibliography{references}

\end{document}
